% ============================================================
% LaTeX Crash Course: Reading main.tex
% ============================================================

% --- Slide 1: What is main.tex? ---
\blueheader
\begin{frame}{What is \texttt{main.tex}?}
\begin{bluebox}{Role of \texttt{main.tex}}{}
  \texttt{main.tex} is the \textbf{root file} of the project. It does two things:
  \begin{itemize}
    \item Sets up the \textbf{preamble}: loads packages, defines colors, macros, and environments.
    \item Pulls in \textbf{content files} with \texttt{\textbackslash include\{...\}}.
  \end{itemize}
\end{bluebox}
\medskip
You \textbf{never write slide content directly in} \texttt{main.tex} --- only setup goes here.
\end{frame}

% --- Slide 2: Document class ---
\blueheader
\begin{frame}{The Document Class}
The very first line sets the document type:
\medskip
\begin{bluefocus}
\texttt{\textbackslash documentclass[11pt, aspectratio=169, handout]\{beamer\}}
\end{bluefocus}
\medskip
\begin{itemize}
  \item \texttt{beamer} --- the presentation class
  \item \texttt{11pt} --- base font size
  \item \texttt{aspectratio=169} --- widescreen 16:9 slides
  \item \texttt{handout} --- strips overlays/animations for printing
\end{itemize}
\end{frame}

% --- Slide 3: Preamble structure ---
\blueheader
\begin{frame}{Preamble: The Setup Zone}
Everything between \texttt{\textbackslash documentclass} and \texttt{\textbackslash begin\{document\}} is the \textbf{preamble}.
\medskip
\begin{bluefocus}
  The preamble in \texttt{main.tex} is organized into sections:
  \begin{enumerate}
    \item Encoding \& typography
    \item Beamer theme \& navigation
    \item Commonly used packages
    \item Hyperlink settings
    \item Custom colors
    \item Frame color shortcuts
    \item Colored box environments
  \end{enumerate}
\end{bluefocus}
\medskip
Each section is marked with a \texttt{\% ---} comment header.
\end{frame}

% --- Slide 4: Packages ---
\blueheader
\begin{frame}{Loading Packages}
Packages extend \LaTeX's capabilities. Syntax:
\medskip
\begin{bluefocus}
\texttt{\textbackslash usepackage[options]\{packagename\}}
\end{bluefocus}
\medskip
Key packages in \texttt{main.tex}:
\begin{itemize}
  \item \texttt{amsmath} --- mathematical typesetting
  \item \texttt{xcolor} --- custom colors
  \item \texttt{tcolorbox} --- colored content boxes
  \item \texttt{fontawesome5} --- icons (\faInfoCircle, \faExclamationTriangle)
  \item \texttt{tikz} / \texttt{pgfplots} --- diagrams and graphs
\end{itemize}
\end{frame}

% --- Slide 5: Custom colors ---
\blueheader
\begin{frame}{Custom Colors}
\texttt{main.tex} defines five named colors used throughout:
\medskip
\begin{bluefocus}
\texttt{\textbackslash colorlet\{myblue\}\{black!35!blue\}}\\
\texttt{\textbackslash colorlet\{myred\}\{black!35!red\}}\\
\texttt{\textbackslash colorlet\{mygreen\}\{black!35!green\}}\\
\texttt{\textbackslash colorlet\{mycyan\}\{cyan\}}\\
\texttt{\textbackslash definecolor\{mymagenta\}\{rgb\}\{0.54,0.0,0.54\}}
\end{bluefocus}
\medskip
\texttt{black!35!blue} means: 35\% black mixed into blue --- giving a \textcolor{myblue}{\textbf{darker, richer blue}}.
\end{frame}

% --- Slide 6: Frame color shortcuts ---
\blueheader
\begin{frame}{Frame Color Shortcuts}
To set the frametitle bar color, call one of these \textbf{before} \texttt{\textbackslash begin\{frame\}}:
\medskip
\begin{bluefocus}
\texttt{\textbackslash blueheader} \quad
\texttt{\textbackslash redheader} \quad
\texttt{\textbackslash greenheader} \quad
\texttt{\textbackslash cyanheader} \quad
\texttt{\textbackslash magentaheader}
\end{bluefocus}
\medskip
Example usage:
\begin{bluefocus}
\texttt{\textbackslash blueheader}\\
\texttt{\textbackslash begin\{frame\}\{My Slide Title\}}\\
\quad \texttt{...content...}\\
\texttt{\textbackslash end\{frame\}}
\end{bluefocus}
\end{frame}

% --- Slide 7: Colored content boxes ---
\blueheader
\begin{frame}{Colored Content Boxes}
Five theorem-style box environments are defined:
\medskip
\begin{itemize}
  \item \texttt{bluebox} --- Definition
  \item \texttt{redbox} --- Theorem
  \item \texttt{greenbox} --- Example
  \item \texttt{cyanbox} --- Exercise
  \item \texttt{magentabox} --- Explore
\end{itemize}
\medskip
Usage requires a \textbf{title} and a \textbf{label}:
\begin{bluefocus}
\texttt{\textbackslash begin\{bluebox\}\{My Title\}\{mylabel\}}\\
\quad \texttt{...content...}\\
\texttt{\textbackslash end\{bluebox\}}
\end{bluefocus}
\end{frame}

% --- Slide 8: Focus boxes ---
\blueheader
\begin{frame}{Focus Boxes}
For \textbf{highlighted text with no title}, use focus boxes:
\medskip
\begin{bluefocus}
\texttt{\textbackslash begin\{bluefocus\}} \quad \texttt{\textbackslash begin\{redfocus\}}\\
\texttt{\textbackslash begin\{greenfocus\}} \quad \texttt{\textbackslash begin\{cyanfocus\}}\\
\texttt{\textbackslash begin\{magentafocus\}}
\end{bluefocus}
\medskip
Example --- this sentence is in a blue focus box:
\begin{bluefocus}
  Focus boxes have a light colored background and no border or title.
\end{bluefocus}
\end{frame}

% --- Slide 9: Warn and info macros ---
\blueheader
\begin{frame}{Warning and Info Macros}
Two icon macros for callouts:
\medskip
\begin{itemize}
  \item \warn
  \item \info
\end{itemize}
\medskip
Custom message versions:
\begin{bluefocus}
\texttt{\textbackslash warnCustomMsg\{Your message here\}}\\
\texttt{\textbackslash infoCustomMsg\{Your message here\}}
\end{bluefocus}
\medskip
\infoCustomMsg{Custom messages let you replace the default text with anything you like.}
\end{frame}

% --- Slide 10: Including content files ---
\blueheader
\begin{frame}{Including Content Files}
Slide content lives in \textbf{separate} \texttt{.tex} files, included at the bottom of \texttt{main.tex}:
\medskip
\begin{bluefocus}
\texttt{\textbackslash include\{Abstract-Linear-Algebra\}}\\
\texttt{\textbackslash include\{French-Pronunciation\}}\\
\texttt{\textbackslash include\{Italian-Pronunciation\}}
\end{bluefocus}
\medskip
\begin{itemize}
  \item Each file contains only frames --- no preamble, no \texttt{\textbackslash begin\{document\}}.
  \item Comment out a line to temporarily exclude that module from the build.
  \item \texttt{\textbackslash include} always starts on a new page; use \texttt{\textbackslash input} if you don't want that.
\end{itemize}
\end{frame}

% --- Slide 11: \newcommand syntax ---
\blueheader
\begin{frame}{\texttt{\textbackslash newcommand} --- Defining Macros}
\texttt{\textbackslash newcommand} creates a reusable shortcut:
\medskip
\begin{bluefocus}
\texttt{\textbackslash newcommand\{\textbackslash name\}[n]\{definition\}}
\end{bluefocus}
\medskip
\begin{itemize}
  \item \texttt{\textbackslash name} --- the new command you are creating
  \item \texttt{[n]} --- \emph{optional}: number of arguments (1--9); omit if none
  \item \texttt{definition} --- what the command expands to; use \texttt{\#1}, \texttt{\#2}\ldots for arguments
\end{itemize}
\medskip
Examples from \texttt{main.tex}:
\begin{bluefocus}
\texttt{\textbackslash newcommand\{\textbackslash warn\}\{\textbackslash textcolor\{orange\}\{\textbackslash textbf\{Not part of the syllabus\}\}\}}\\[4pt]
\texttt{\textbackslash newcommand\{\textbackslash warnCustomMsg\}[2][orange]\{\textbackslash textcolor\{\#1\}\{\textbackslash textbf\{\#2\}\}\}}
\end{bluefocus}
\medskip
The second example uses \texttt{[2][orange]}: 2 arguments, first has default value \texttt{orange}.
\end{frame}

% --- Slide 12: \tcbset syntax ---
\blueheader
\begin{frame}{\texttt{\textbackslash tcbset} --- Global Box Defaults}
\texttt{\textbackslash tcbset} sets \textbf{default options} for all \texttt{tcolorbox} environments that follow:
\medskip
\begin{bluefocus}
\texttt{\textbackslash tcbset\{key=value, key=value, ...\}}
\end{bluefocus}
\medskip
From \texttt{main.tex}:
\begin{bluefocus}
\texttt{\textbackslash tcbset\{}\\
\quad \texttt{colback=white,} \hfill \textit{background color}\\
\quad \texttt{fonttitle=\textbackslash bfseries,} \hfill \textit{bold titles}\\
\quad \texttt{breakable,} \hfill \textit{allow page breaks}\\
\quad \texttt{enhanced,} \hfill \textit{enable advanced drawing}\\
\quad \texttt{top=2pt, bottom=0pt, left=0pt, right=6pt,}\\
\quad \texttt{boxsep=4pt}\\
\texttt{\}}
\end{bluefocus}
\medskip
Any \texttt{tcolorbox} or \texttt{\textbackslash newtcbtheorem} box inherits these defaults automatically.
\end{frame}

% --- Slide 13: \newtcbtheorem syntax ---
\blueheader
\begin{frame}{\texttt{\textbackslash newtcbtheorem} --- Numbered Box Environments}
\texttt{\textbackslash newtcbtheorem} creates a \textbf{numbered, titled box environment}:
\medskip
\begin{bluefocus}
\texttt{\textbackslash newtcbtheorem[counter options]\{envname\}\{Title\}\{tcolorbox options\}\{prefix\}}
\end{bluefocus}
\medskip
\begin{itemize}
  \item \texttt{[counter options]} --- e.g.\ \texttt{number within=section} or \texttt{use counter from=bluebox}
  \item \texttt{envname} --- the environment name you use in your document
  \item \texttt{Title} --- the label printed in the box header (e.g.\ \emph{Definition}, \emph{Theorem})
  \item \texttt{\{tcolorbox options\}} --- extra styling, e.g.\ \texttt{colframe=myblue}
  \item \texttt{prefix} --- namespace for \texttt{\textbackslash ref} labels, e.g.\ \texttt{bluebox}
\end{itemize}
\medskip
\begin{bluefocus}
\texttt{\textbackslash newtcbtheorem\{bluebox\}\{Definition\}\{colframe=myblue\}\{bluebox\}}\\[2pt]
Then use: \texttt{\textbackslash begin\{bluebox\}\{My Title\}\{mylabel\} ... \textbackslash end\{bluebox\}}
\end{bluefocus}
\end{frame}
